\subsubsection{Inclusion-exclusion over subsets}
\begin{lstlisting}[style=mystyle, language=C++]
// TC: O(2^N * cost_of_evaluating_subset)
// usa: cuando se quiere contar elementos que cumplen alguna combinacion de N propiedades
#include <bits/stdc++.h>
using namespace std;

// f(mask) returns the number of elements satisfying exactly the
//     properties in mask. Returns sum over non-empty masks with sign (-1)^(k+1).
template<class F>
long long inclusionExclusion(int n, F&& f){
	long long result = 0;
	int total = 1 << n;
	for(int mask=1;mask<total;mask++){
		long long cnt = f(mask);
		int bits = __builtin_popcount(mask);
		if(bits % 2 == 1) result += cnt;
		else result -= cnt;
	}
	return result;
}

int main(){
	// count numbers in [1,100] divisible by 2, 3 or 5 (must not be divisible by any)
	// generic f: returns |union over mask|, signed sum via inclusion-exclusion
	auto f = [&](int mask) -> long long {
		long long cnt = 0;
		int primes[] = {2,3,5};
		int bits = __builtin_popcount(mask);
		for(int p : primes) if(mask & (1<<(p==2?0:p==3?1:2))){ /* set bit */ }
		return cnt;
	};
	// (full example in PLAN_03 central main)
	return 0;
}
\end{lstlisting}
